% Section 11.3, practice test 2, translated from sv/chapters/11/dugga.tex.

\needspace{16\baselineskip}
\section{Practice test 2}
\label{c11:sec:ovningsdugga}
The practice test covers Chapters~\ref{c7:ch}--\ref{c10:ch}: vectors, functions, trigonometry,
logarithms and decibels. It is set out like a real test, so take it on your own and in one sitting
before you look in the answer key, where the answers are.

\begin{aside}
\textbf{How the test works.} Permitted aids are writing materials, any calculator and the formulas
needed from the course's formula sheet. The test has five questions worth $1.0$ point each, so at
most $5.0$ points. Every question requires a worked solution with the answer stated, including the
unit where needed. The real test is set out the same way: it does not affect your grade, but at
least $3.5$ of $5.0$ points gives $0.5$ bonus points on the regular exam.
\end{aside}

\prov{od2}

\provuppgift{1}{1.0}
You have two vectors $\mathbf{u} = (1, 2)$ and $\mathbf{v} = (-3, 4)$.
\begin{parts}
  \item Draw the vectors in a coordinate system.
  \item Calculate the magnitudes $\abs{\mathbf{u}}$ and $\abs{\mathbf{v}}$ of the vectors.
  \item Calculate the angles $v_u$ and $v_v$ of the vectors.
  \item Calculate a third vector $\mathbf{w} = 2\mathbf{u} - \mathbf{v}$.
\end{parts}

\provuppgift{2}{1.0}
Simplify the following rational expression and state which values of $x$ are not permitted:
\[
  f(x) = \frac{x^2 - x - 2}{x - 2}
\]

\provuppgift{3}{1.0}
When a capacitor is discharged through a resistor, the voltage decreases exponentially with time.
The voltage $u(t)$ can be described by
\[
  u(t) = U_0 e^{-t/(RC)},
\]
where
\begin{itemize}
  \item $u(t)$ is the voltage across the capacitor at the time $t$,
  \item $U_0$ is the initial voltage in V,
  \item $RC$ is the circuit's time constant in seconds,
  \item $t$ is the time in seconds.
\end{itemize}
A capacitor with the capacitance $470\,\mu\text{F}$ is connected to a $10\,\text{k}\Omega$
resistor. It is initially charged to $12\,\text{V}$ and starts to discharge at $t = 0$.
\begin{parts}
  \item Calculate the voltage after five seconds.
  \item Find the function's domain and range for the time interval
    $0 \leq t \leq 20\,\text{s}$.
  \item Calculate after how long the voltage has fallen to half its original value.
\end{parts}

\provuppgift{4}{1.0}
An AC voltage has the amplitude $5\,\text{V}$, the frequency $100\,\text{Hz}$ and the phase
$90^{\circ}$.
\begin{parts}
  \item Find the equation $u(t)$ of the AC voltage. Give the phase in radians.
  \item Draw the sine curve of the AC voltage over one period $T$.
\end{parts}

\provuppgift{5}{1.0}
An audio amplifier has a gain of $32\,\text{dB}$. By what factor does it amplify the voltage?
